\chapter{Towards the Synthesis of a Self-Healing Polymer} % enter the name of the chapter here
\label{cha:chem} % enter the chapter label here (for cross referencing)

This study aims to lay initial ground-work for the development of a self-healing polymer. This involves Kapton\textsuperscript{\textregistered}, a commercially available polyimide, well known for highly resilient properties. This section covers the initial process for the alteration of polyimide via hydrolysis in order to expose reactive carboxylic acid (COOH) groups and the subsequent reactions with these exposed COOH groups. Hydrolysis is a process in which polyimide is broken down into a precursor, polyamic acid, this exposes regions in which additional reactivity can be added. One route to the development of a self healing polymer is through the formation of crosslinks at these exposed COOH sites. Specific types of crosslinks can dynamically break and reform to achieve a material that can repair itself. A proof of concept reaction between polyamic acid and ethylene carbonate, a crosslinking reagent, was attempted as a preliminary demonstration.


\section{Results}
\subsection{Hydrolysis of the Polyimide}
The hydrolysis of the polyimide follows a standard mechanism for hydrolysis of an amide to a carboxylic acid using sodium hydroxide (NaOH). These show the hydroxy group attacking the carbonyl which then breaks the C-N bond and gives the carboxylic acid free.

During the formation of the hydrolysed version, the polymer sheets were converted into a light orange suspension which was filtered to give a pale orange precipitate.

\subsubsection{Structural Analysis}
Structural determination was conducted via Infrared spectroscopy. A comparison was made between the attempted hydrolysis of the sheets and a purchased form (polyamic acid). 

\subsection{Cross-linking of the Polyimide}
Cross-linking attempts follow a mechanism as proposed by ( put reference here), where ethylene carbonate is used to form an ester with the carboxylic acids on neighbouring polymer chains, thereby joining them together with a "cross-link". This was performed with three different salts ( TEAB, $ZnCl_2$, $FeCl_3$) as catalysts. 

During the cross-linking reaction with TEAB, an orange-brown suspension was formed that when washed, filtered and freeze dried gave a pale orange powder. With $ZnCl_2$, the liquid was a dark orange clear liquid that, when water was added, gave a single clump of grey polymeric residue. This was freeze dried to remove the excess water and then analysed. The cross-linking reaction with the $FeCl_3$, gave a dark orange/brown gel, which gave a grey gum-like texture in small pieces when water was added. This was also frozen and analysed. 

\subsubsection{Infrared Spectrum analysis}



% \section{Hydrolysis of a Polyimide}
% The hydrolysis of the poylimide (kapton) is needed to open an amide bond on the imide ring to reveal a carrboxylic acid group whilst the bolymer back bone is still intact. This will make the polymer soluble, however since the polymer is insoluble currently the only attack will be on the polymer surface. This was completed with sodium hydroxide in (solvent) and in the monowave reactor for (x time)

% The structure of the proposed polyimide is shown in figure (x). This structure has multiple imide rings that we hypothesise are a good starting point for introducing reactivity. 

% The partial hydrolysis of these rings was performed with Sodium hydroxide solution (in dmso). This reaction was performed in the monowave reactor. The monowave reactor is a reaction vessel that is within a machine that can reduce the pressure on the environment. This is useful as the boiling point of the solvent is increase with lower pressure. Therefore a reaction that can only be performed above a certain temperature like this one can be done in solvent with a lower boiling point than the reactuion requirements.

% \subsection{}


% %input graph about changinh boiling point of solvent with pressure

% \subsection{Comparison of Hydrolysed Polyimide versus Purchased Polyamic Acid}
% the hydrolysed polyimide should have a similar ir as a purchased form of the polyimide as both have a exposed carboxylic acid group, therefore there will be a characteristic shift at (wavenumber).


% \subsection{Spectral Analysis}
% Structural determination was performed with infrared (IR) spectroscopy and Ultraviolet - visible (UV-Vis) spectroscopy. Key IR peaks for the unmodified Kapton are (x), the key peask for the modified kapton are (x) and the key peaks for the purchased Polyamic acid are (y). 



% \section{Crosslinking of Polyamic Acid}

% crosslinking of polyamic acid was done in a reaction with the ethylene carbonate and an ammonium bromide salt this gave the reaction as in a similar process as in (this paper ). the carboxylic acid groups react with the ethylene carbonate to form either an ester with a hydroxy after an alkyl chain on one of the polymer chains or it will like the polymer chains together by forming a diester. 


% Using the monowave uses a 10mL reaction tube that was filled with 4ml of reagents to leave room for the co2 that is released in the reaction. The volume of co2 released was calculated using
% $PV = nRT$ which was rearranged to $n = \frac{PV}{RT}$. The monowave 50 has a maximum pressure of 20 bar, so we opted for a maximum of 5 bar for safety. this was combined with a temperature of 150$\degree$C to give: 

% \[n = \frac{PV}{RT}\]
% \[n = \]


% THis was the total number of co2 mo

% \section{Addition of Diels-Alder Groups}

% \section{Experimental}
% Kapton (poly-oxydiphenylene-pyromellitimide) sheets were cut into small pieces ($\approx$ 2mm$^2$). These were added to a solution of NaOH (vol and mol here) and heated to 120 $\degree$C in a monowave reactor set to ramp for 20 mins and hold for 20 and 40 mins successively. The reaction gave no observable difference in the Kapton sheet

\section{Experimental}
\subsection{The Hydrolysis of a Polyimide}
Poly..... (Kapton) film (0.2g, thickness of x) was cut into pieces of approx. 3mm$^2$, these were mixed with a solution of NaOH () in a reaction tube. This tube was placed into a monowave reactor that was set to ramp for 20 minutes and hold for 20 minutes at 100 $\degree$C. A sample was taken and the remaining mixture was reacted again on the monowave for ramp 20 mins and hold 20 minutes at 100$\degree$C. There was no observable difference in the Kapton film. Then this was thoroughly washed with water and IR was collected.

\subsection{Cross-linking with TEAB and Ethylene carbonate}
A mixture of polyamic acid (0.51mL, conc) was mixed with ethylene carbonate () and tetraethylammonium bromide (). This was mixed in a monowave reactor set to ramp for 20 minutes and hold for 1 hour at 150 $\degree$C. This gave an orange suspension that was dried to give a pale orange powder as the final product (grams, yield).
IR (solid): 

\subsection{Cross-linking with $ZnCl_2$ and Ethylene carbonate}
A mixture of polyamic acid (0.51mL, conc) was mixed with ethylene carbonate () and Zinc chloride (). This was mixed in a monowave reactor set to ramp for 20 minutes and hold for 1 hour at 100 $\degree$C. This gave dark orange clear liquid that was dried to give a () as the final product (grams, yield).
IR (solid): 

\subsection{Cross-linking with $FeCl_3$ and Ethylene carbonate}
A mixture of polyamic acid (0.51mL, conc) was mixed with ethylene carbonate () and Iron(III) chloride (). This was mixed in a monowave reactor set to ramp for 20 minutes and hold for 1 hour at 100 $\degree$C. This gave a dark orange clear gel that was dried to give a () as the final product (grams, yield).
IR (solid): 